% ID: SEL-D0005-OA-E03
% ESTADO: cotejada. Segunda pasada visual completa y borrador compilado revisados.
% Fuente: Selectividad/sel_2019_matematicas_aplicadas/Reserva A/Exámenes Mat Aplicadas CCSS - A y B.pdf
% SHA-256: d28871f7f203cf1be9e479e9d642cdb6b1623704a8b76a2b57cfa2d12950f293
% Página: 1; opción A; ejercicio 3.

Una cooperativa envasa zumos de naranja, zumos de piña y zumos de melocotón
en botellas de 1 litro y de 2 litros. Se sabe que el 60\,\% de las botellas
son de zumo de naranja y el 30\,\% de piña. Además, el 80\,\% de las
botellas de zumo de naranja y el 70\,\% de las de zumo de piña son de
2 litros, mientras que el 60\,\% de las de melocotón son botellas de 1 litro.
Se elige al azar una botella envasada por la cooperativa.
\begin{enumerate}[label=\alph*)]
  \item \textbf{(1 punto)} Calcule la probabilidad de que la botella sea de
  2 litros.
  \item \textbf{(0.75 puntos)} Calcule la probabilidad de que el zumo sea de
  naranja, sabiendo que la botella es de 2 litros.
  \item \textbf{(0.75 puntos)} Calcule la probabilidad de que el zumo sea de
  melocotón, sabiendo que la botella es de 1 litro.
\end{enumerate}
